\documentclass[11pt]{article}

\usepackage[margin=1in]{geometry}
\usepackage{setspace}
\usepackage{hyperref}
\usepackage{titlesec}
\usepackage{amsmath}
\usepackage{amssymb}
\usepackage{booktabs}
\usepackage{multirow}

\titleformat{\section}{\large\bfseries}{\thesection}{1em}{}
\titleformat{\subsection}{\normalsize\bfseries}{\thesubsection}{1em}{}

\title{\textbf{LLM-Augmented Multi-Agent Evacuation Simulation\\
Under Dynamic Disaster Conditions}\\[0.4em]
\large CS 6960 Capstone -- Final Report}
\author{Sheng Hung Tseng}
\date{}

\begin{document}

\maketitle

\onehalfspacing

% ─────────────────────────────────────────────
\section{Introduction and Motivation}

Campus mass evacuation under dynamic disaster conditions presents a challenging planning problem: the road network degrades in real time, individual evacuees have heterogeneous mobility and behavioral characteristics, and no single agent has access to the full global state. Standard heuristic policies such as nearest-shelter routing are distance-aware but fail to adapt to mid-evacuation route failures~\cite{chen2018}, while prior RL-based approaches~\cite{mnih2015} typically assume a homogeneous agent population.

This project designs and evaluates a three-layer decision-support system for campus evacuation on the University of Utah road network. The system combines LLM-based behavioral profiling, zone-level coordination, and a persona-aware DRQN navigation policy trained on LLM-calibrated disaster parameters. The central question is whether integrating population heterogeneity into both simulation design and policy training leads to more robust and equitable evacuation outcomes.

% ─────────────────────────────────────────────
\section{System Architecture}

The system is organized into three layers that operate at increasing levels of granularity.

\textbf{Layer 1 -- LLM Behavioral Profiling.} Llama 3.3 70B~\cite{meta2024llama3} (via Groq API) generates 20 heterogeneous persona profiles across four campus roles: students (60\%), faculty (15\%), staff (20\%), and visitors (5\%). Each persona is characterized by six quantitative behavioral fields:

\begin{itemize}
  \item \textbf{walk\_speed\_multiplier}: Relative walking speed compared to a baseline pedestrian (1.0 = normal). Values above 1.0 represent faster movers (e.g., athletes); below 1.0 represent slower movers (e.g., mobility-impaired individuals).
  \item \textbf{compliance\_rate}: Probability that the agent follows official evacuation instructions (e.g., takes the recommended route rather than self-routing). Higher compliance reflects greater trust in authority and familiarity with emergency protocols.
  \item \textbf{panic\_level}: Degree of psychological stress during the disaster event (0 = calm, 1 = fully panicked). Panic degrades both decision-making quality and navigation accuracy.
  \item \textbf{observation\_error\_multiplier}: Scaling factor applied to the agent's perception noise. A value of 1.0 means standard noise; higher values reflect impaired situational awareness due to stress, unfamiliarity, or physical limitations.
  \item \textbf{decision\_delay\_steps}: Number of simulation steps the agent waits before beginning to move after the evacuation signal, modeling hesitation or slow reaction.
  \item \textbf{shelter\_familiarity}: Prior knowledge of shelter locations on campus (0 = completely unfamiliar, 1 = fully familiar). Low familiarity leads to longer search times and increased likelihood of choosing suboptimal routes.
\end{itemize}

Panic level modulates two fields dynamically at inference time, reflecting that heightened stress simultaneously reduces compliance and worsens perception:
\[
  \text{eff\_compliance} = \text{compliance} \times (1 - 0.5 \times \text{panic}), \quad
  \text{eff\_obs\_error} = \text{obs\_error} \times (1 + \text{panic})
\]

\textbf{Layer 2 -- Zone-Level Coordination.} A zone assignment module partitions the campus into geographic zones and maps each zone to primary and backup shelters subject to capacity constraints. An LLM-assisted coordinator (also via Groq) interprets aggregated zone status and generates natural-language route recommendations following a ReAct-style~\cite{yao2023react} reasoning-then-action pipeline: the model first reasons about current hazard conditions and shelter availability, then produces a concrete evacuation directive. A REST API (FastAPI) exposes this layer for interactive querying.

\textbf{Layer 3 -- Persona-Aware DRQN Navigation.} A Deep Recurrent Q-Network with GRU memory and Dueling DQN heads navigates individual agents on the OSM pedestrian graph. The observation vector encodes both environmental state and persona features, enabling the policy to adapt routing decisions to agent-specific behavioral characteristics.

% ─────────────────────────────────────────────
\section{Methodology}

\subsection{Graph Environment}

The campus road network is extracted from OpenStreetMap~\cite{osm} using OSMnx~\cite{osmnx} and represented as a directed graph $G = (V, E)$. Each edge $e = (u,v) \in E$ carries a dynamic weight:
\[
  w(e, t) = d(u,v) \cdot \bigl(1 + \alpha_{\mathrm{slope}} \cdot \mathrm{slope}(e) + \alpha_{\mathrm{snow}} \cdot h_{\mathrm{snow}}(e, t)\bigr)
\]
where $d(u,v)$ is Euclidean distance, $\mathrm{slope}(e)$ is derived from USGS elevation data, and $h_{\mathrm{snow}}(e,t)$ is the current snow/debris depth. When $h_{\mathrm{snow}}(e,t)$ exceeds a blockage threshold $\tau$, the edge is stochastically removed from the routing graph.

\subsection{Hazard Processes}

Three disaster types are modeled.

\textbf{Blizzard.} Snow depth accumulates stochastically:
\[
  h_{\mathrm{snow}}(e, t+1) = h_{\mathrm{snow}}(e, t) + \epsilon_t, \quad \epsilon_t \sim \mathcal{N}(\mu_s, \sigma_s^2)
\]
Roads degrade progressively over time with minimal initial damage.

\textbf{Earthquake.} Road segments are blocked at simulation onset with probability $p_{\mathrm{init}}$ (structural collapse), then continue to fail stochastically each step (aftershocks). Blockage is front-loaded rather than gradual.

\textbf{Compound Disaster.} Both processes operate simultaneously: immediate structural damage (earthquake) combined with progressive snow accumulation (blizzard). Initial blockage is the maximum of both components at equivalent severity.

\textbf{LLM-calibrated parameters.} Rather than hand-tuning scenario parameters, Llama 3.3 70B~\cite{meta2024llama3} was prompted with physical text descriptions of each parameter and asked to produce values consistent with its knowledge of real-world disasters. The LLM output was post-processed with safety clamps and saved to \texttt{scenarios/llm\_severity\_presets.json}, which the simulation loads automatically. Table~\ref{tab:llm_params} summarizes the LLM-generated blizzard parameters.

\begin{table}[h]
\centering
\caption{LLM-generated blizzard severity parameters (key fields).}
\label{tab:llm_params}
\begin{tabular}{lcccc}
\toprule
\textbf{Severity} & $p_{\mathrm{init}}$ & $\tau$ (threshold) & $p_{\mathrm{snow\rightarrow block}}$ & $\alpha_{\mathrm{snow}}$ \\
\midrule
Light    & 0.000 & 0.80 & 0.010 & 1.2 \\
Moderate & 0.010 & 0.40 & 0.020 & 1.5 \\
Severe   & 0.010 & 0.60 & 0.030 & 1.8 \\
Extreme  & 0.050 & 0.40 & 0.030 & 2.5 \\
\bottomrule
\end{tabular}
\end{table}

\subsection{Persona-Aware Observation Space}

Each agent's observation vector $o_t \in \mathbb{R}^{42}$ consists of three components:

\begin{enumerate}
  \item \textbf{Spatial features} (7 dims): displacement to current sub-goal $(dx, dy)$, displacement to final shelter $(dx', dy')$, local blocked-edge ratio, local mean snow depth, normalized timestep $t / T$.
  \item \textbf{Persona features} (5 dims): normalized speed ($v / 1.5$), panic level, shelter familiarity, effective compliance, effective observation error.
  \item \textbf{Neighbor features} (5 dims $\times$ 6 neighbors = 30 dims): for each of the top-$k$ candidate neighbors -- existence flag, graph progress, normalized edge length, edge snow depth, edge blocked flag.
\end{enumerate}

The persona feature block allows the policy to modulate routing aggressiveness based on agent type, without requiring separate models per persona.

\subsection{DRQN Architecture}

The navigation policy is implemented following the DRQN framework~\cite{hausknecht2015} with a Dueling DQN~\cite{wang2016} head and GRU~\cite{cho2014gru} memory in place of LSTM:
\begin{itemize}
  \item \textbf{Encoder:} Fully connected layer mapping $o_t \in \mathbb{R}^{42}$ to 128-dim embedding (ReLU).
  \item \textbf{GRU:} Single-layer GRU with hidden size 128, maintaining a recurrent belief state across steps.
  \item \textbf{Dueling head:} Separate value stream $V(s)$ and advantage stream $A(s,a)$; combined as $Q(s,a) = V(s) + A(s,a) - \bar{A}(s)$.
  \item \textbf{Action space:} Stay in place or move to one of $k=6$ top-ranked candidate neighbors.
\end{itemize}

Training uses Double DQN~\cite{vanHasselt2016} target updates, a sequence replay buffer (episode fragments of length 8), and a burn-in period of 4 steps. A 4-stage curriculum~\cite{bengio2009curriculum} progressively increases disaster severity (light $\rightarrow$ moderate $\rightarrow$ severe $\rightarrow$ extreme), with persona randomly sampled each episode from all 20 profiles.

% ─────────────────────────────────────────────
\section{Experimental Setup}

\textbf{Population.} Each simulation run deploys 40 pedestrian agents and 15 vehicles, reflecting the default campus scenario configuration. Agents are sampled from the 20 persona types according to campus role proportions (student 60\%, staff 20\%, faculty 15\%, visitor 5\%).

\textbf{Scenarios.} Three disaster types (blizzard, earthquake, compound) $\times$ four severity levels (light, moderate, severe, extreme) = 12 scenario conditions. Each condition is evaluated over 20 independent runs with different random seeds.

\textbf{Baselines.} Round Robin (naive assignment) and Nearest Shelter (distance-aware greedy) serve as heuristic references. The DRQN policy is evaluated as the primary learned policy.

\textbf{Evaluation metrics:}
\begin{itemize}
  \item \textbf{Reached Rate}: fraction of agents successfully reaching a shelter.
  \item \textbf{Avg.\ Exposure}: total accumulated hazard exposure per agent.
  \item \textbf{Per-Persona Reached Rate}: reached rate broken down by all 20 persona types.
  \item \textbf{Fairness Gap}: max persona reached rate minus min persona reached rate.
\end{itemize}

% ─────────────────────────────────────────────
\section{Results}

\subsection{Baseline vs.\ DRQN under Blizzard}

Table~\ref{tab:baseline_vs_drqn} compares heuristic policies and DRQN under blizzard conditions across severity levels (40 agents, 20 runs each). Round Robin and Nearest Shelter numbers are from the baseline scenario evaluation; DRQN results are from the blizzard severity sweep.

\begin{table}[h]
\centering
\caption{Reached rate comparison: baselines vs.\ DRQN under blizzard (40 agents, 20 runs).}
\label{tab:baseline_vs_drqn}
\begin{tabular}{lcccc}
\toprule
\textbf{Scenario} & \textbf{Round Robin} & \textbf{Nearest} & \textbf{DRQN (prev.)} & \textbf{DRQN (persona-aware)} \\
\midrule
Baseline (no disaster) & 0.265 & 0.291 & 0.963 & -- \\
Blizzard Light    & --    & --    & 0.845 & \textbf{0.790} \\
Blizzard Moderate & --    & --    & 0.752 & \textbf{0.724} \\
Blizzard Severe   & --    & --    & 0.640 & \textbf{0.678} \\
Blizzard Extreme  & --    & --    & 0.519 & \textbf{0.545} \\
\bottomrule
\end{tabular}
\end{table}

DRQN outperforms both heuristics by a large margin under baseline conditions (+262\% over round robin). Under blizzard conditions, the persona-aware model achieves comparable or improved performance relative to the previous DRQN model despite a harder training distribution from LLM-calibrated parameters. The slight gain at extreme severity (0.545 vs.\ 0.519) suggests that exposure to diverse persona dynamics during training aids generalization under the most challenging conditions.

\subsection{DRQN Across All Disaster Types and Severities}

Table~\ref{tab:sweep} summarizes DRQN performance across all 12 scenario conditions using the 200-agent population.

\begin{table}[h]
\centering
\caption{DRQN reached rate and avg.\ exposure across disaster types and severities (40 agents, 20 runs, persona-aware model).}
\label{tab:sweep}
\begin{tabular}{llcc}
\toprule
\textbf{Disaster} & \textbf{Severity} & \textbf{Reached Rate} & \textbf{Avg.\ Exposure} \\
\midrule
\multirow{4}{*}{Blizzard}
  & Light    & 0.790 & 31.3 \\
  & Moderate & 0.724 & 70.7 \\
  & Severe   & 0.678 & 84.4 \\
  & Extreme  & 0.545 & 72.5 \\
\midrule
\multirow{4}{*}{Earthquake}
  & Light    & 0.384 & 112.9 \\
  & Moderate & 0.396 & 72.7 \\
  & Severe   & 0.404 & 62.1 \\
  & Extreme  & 0.449 & 27.3 \\
\midrule
\multirow{4}{*}{Compound}
  & Light    & 0.625 & 76.4 \\
  & Moderate & 0.374 & 65.5 \\
  & Severe   & 0.363 & 63.7 \\
  & Extreme  & 0.417 & 37.7 \\
\bottomrule
\end{tabular}
\end{table}

Blizzard yields the highest reached rates because road degradation is gradual, giving the policy time to reroute. Earthquake produces lower rates and higher exposure due to sudden mid-route failures. Compound disaster shows the steepest decline under moderate severity, where both hazard processes compound simultaneously; the slight recovery at extreme severity may reflect that heavily damaged networks force shorter, faster routes.

\subsection{Per-Persona Fairness Analysis}

A key contribution of this work is the analysis of evacuation equity across the 20 persona types. Table~\ref{tab:fairness} reports per-role reached rates under blizzard conditions.

\begin{table}[h]
\centering
\caption{Per-role reached rate under blizzard (40 agents, persona-aware DRQN).}
\label{tab:fairness}
\begin{tabular}{lcccc}
\toprule
\textbf{Severity} & \textbf{Student} & \textbf{Faculty} & \textbf{Staff} & \textbf{Visitor} \\
\midrule
Light    & 0.751 & 0.829 & 0.847 & 0.672 \\
Moderate & 0.684 & 0.760 & 0.831 & 0.544 \\
Severe   & 0.598 & 0.668 & 0.705 & 0.477 \\
Extreme  & 0.553 & 0.666 & 0.713 & 0.410 \\
\bottomrule
\end{tabular}
\end{table}

Staff consistently achieves the highest reached rate due to high campus familiarity, high compliance, and low panic. Visitors are the most vulnerable group across all conditions, with the gap widening under severe and extreme scenarios.

At the persona level, the fairness gap at extreme blizzard severity is 0.708 (campus\_security: 0.833 vs.\ conference\_attendee: 0.125). The five most vulnerable personas at extreme severity are:

\begin{enumerate}
  \item \textbf{conference\_attendee} (visitor): 0.125 -- low familiarity, high observation error
  \item \textbf{prospective\_student\_with\_parent} (visitor): 0.250 -- low familiarity, moderate panic
  \item \textbf{student\_with\_anxiety} (student): 0.384 -- panic-induced compliance collapse
  \item \textbf{adjunct\_instructor} (faculty): 0.420 -- low familiarity vs.\ senior faculty
  \item \textbf{mobility\_impaired} (visitor): 0.429 -- speed bottleneck under severe conditions
\end{enumerate}

These results show that panic-level modulation ($\text{eff\_compliance} = \text{compliance} \times (1 - 0.5 \times \text{panic})$) is the primary driver of poor performance in anxiety-prone personas, while low shelter familiarity dominates for visitor-type agents unfamiliar with the campus layout.

\subsection{Cross-Disaster Fairness Analysis}

Table~\ref{tab:cross_fairness} reports per-role reached rates at extreme severity across all three disaster types, revealing how evacuation equity shifts with disaster character.

\begin{table}[h]
\centering
\caption{Per-role reached rate at extreme severity across disaster types (40 agents, persona-aware DRQN).}
\label{tab:cross_fairness}
\begin{tabular}{lcccc}
\toprule
\textbf{Disaster} & \textbf{Student} & \textbf{Faculty} & \textbf{Staff} & \textbf{Visitor} \\
\midrule
Blizzard   & 0.522 & 0.579 & 0.557 & 0.606 \\
Earthquake & 0.334 & 0.543 & \textbf{0.796} & 0.136 \\
Compound   & 0.308 & 0.535 & \textbf{0.718} & 0.170 \\
\bottomrule
\end{tabular}
\end{table}

The disaster type fundamentally reshapes the fairness landscape. Under earthquake and compound conditions, staff achieve the highest reached rates (0.796 and 0.718) while visitors are the most vulnerable (0.136 and 0.170), an asymmetry not present in blizzard. This is driven by the obs\_error and familiarity fields: staff personas (campus\_security, facilities\_staff) have obs\_error $= 0.50$ and familiarity $= 0.95$--$1.00$, enabling them to navigate effectively even when 80\% of roads are initially blocked (earthquake extreme $p_{\mathrm{init}} = 0.80$). Visitor personas with obs\_error $> 2.50$ and familiarity $= 0.10$--$0.20$ fail catastrophically under the same conditions.

The maximum per-persona fairness gap peaks at earthquake extreme (campus\_security: 1.000 vs.\ mobility\_impaired: 0.000, gap = 1.000), followed by compound extreme (campus\_security: 0.889 vs.\ conference\_attendee: 0.125, gap = 0.764), and blizzard extreme (gap = 0.557). The mobility\_impaired persona achieves zero reached rate under earthquake extreme because its low walk speed (0.30$\times$) combined with panic-amplified observation error leaves it unable to identify any unblocked exit path when initial blockage is 80\%.

% ─────────────────────────────────────────────
\section{Analysis and Discussion}

\textbf{LLM parameter calibration.} Compared to hand-tuned parameters, the LLM-generated values assigned notably lower blockage thresholds ($\tau$) for moderate and extreme conditions, reflecting a physical judgment that heavily accumulated snow converts to road closures more aggressively than the original conservative values. This produces a harder training distribution, which may explain the slight decrease in extreme blizzard performance of the persona-aware model.

\textbf{Persona features in the obs vector.} Including the 5-dimensional persona feature block in the DRQN observation appears to improve generalization across moderate and severe conditions. The policy can now distinguish a high-panic freshman (who requires more conservative routing) from a calm campus security agent (who can tolerate riskier shortcuts), rather than treating all agents identically.

\textbf{Earthquake non-monotonicity.} Under the LLM-calibrated earthquake parameters, reached rate increases monotonically from light to extreme (0.384 → 0.396 → 0.404 → 0.449), which is counter-intuitive. This pattern is consistent with the LLM-generated parameters assigning very high initial blockage at extreme severity ($p_{\mathrm{init}} = 0.80$): when most long routes are immediately cut off, agents are forced onto shorter, more reliable paths, reducing failure from mid-route blockage. Average exposure also drops sharply (112.9 → 27.3), confirming that extreme earthquakes produce faster but lower-mileage evacuations.

\textbf{Limitations.} Several limitations constrain the current evaluation. The 40-agent population does not model shelter capacity queuing or congestion effects. The fairness analysis is statistically limited for low-frequency personas (visitor roles average 1--2 agents per run), flagged with confidence warnings in the output. The DRQN policy is trained only on blizzard conditions and applied zero-shot to earthquake and compound scenarios without disaster-specific fine-tuning, which may explain the lower absolute performance on those disaster types.

% ─────────────────────────────────────────────
\section{Implementation Overview}

The project is implemented in Python and consists of the following components:

\begin{itemize}
  \item \textbf{evac\_env.py}: OSM-based evacuation environment with dynamic snow/blockage.
  \item \textbf{drqn\_minimal.py}: Persona-aware DRQN training (42-dim obs, GRU + Dueling DQN, curriculum learning, persona sampling per episode).
  \item \textbf{batch\_runner.py}: Multi-agent simulation runner; wires actual agent persona fields into obs at inference time.
  \item \textbf{llm\_behavior\_profiler.py}: Layer 1 LLM persona generation via Groq.
  \item \textbf{llm\_zone\_coordinator.py}: Layer 2 zone-level LLM coordination.
  \item \textbf{llm\_scenario\_generator.py}: LLM-driven disaster parameter generation; outputs \texttt{scenarios/llm\_severity\_presets.json}.
  \item \textbf{scenario\_loader.py}: Loads LLM presets at import time; falls back to hardcoded tables.
  \item \textbf{advisor\_api.py}: FastAPI REST service exposing full three-layer pipeline.
  \item \textbf{analyze\_persona\_fairness.py}: Per-persona equity analysis with confidence flagging.
  \item \textbf{run\_disaster\_severity\_sweep.py}: Automated 3-disaster $\times$ 4-severity sweep runner.
\end{itemize}

% ─────────────────────────────────────────────
\section{Conclusion}

This project demonstrates that integrating LLM-based behavioral heterogeneity into both the simulation design and policy training pipeline leads to measurable improvements in evacuation outcome analysis. The persona-aware DRQN substantially outperforms heuristic baselines across all disaster conditions, achieving reached rates of 0.79 under light blizzard and maintaining above 0.54 under extreme conditions -- compared to near-zero performance for both heuristic policies. The fairness analysis reveals significant equity gaps between population groups (gap of 0.708 at extreme blizzard severity), motivating future work on fairness-constrained policy training that explicitly accounts for vulnerable persona characteristics.

% ─────────────────────────────────────────────
\begin{thebibliography}{9}

\bibitem{hausknecht2015}
M.~Hausknecht and P.~Stone,
``Deep recurrent Q-learning for partially observable MDPs,''
\textit{arXiv preprint arXiv:1507.06527}, 2015.

\bibitem{wang2016}
Z.~Wang, T.~Schaul, M.~Hessel, H.~van Hasselt, M.~Lanctot, and N.~de Freitas,
``Dueling network architectures for deep reinforcement learning,''
in \textit{Proceedings of the 33rd International Conference on Machine Learning (ICML)}, pp.~1995--2003, 2016.

\bibitem{vanHasselt2016}
H.~van Hasselt, A.~Guez, and D.~Silver,
``Deep reinforcement learning with double Q-learning,''
in \textit{Proceedings of the 30th AAAI Conference on Artificial Intelligence}, pp.~2094--2100, 2016.

\bibitem{osmnx}
G.~Boeing,
``OSMnx: New methods for acquiring, constructing, analyzing, and visualizing complex street networks,''
\textit{Computers, Environment and Urban Systems}, vol.~65, pp.~126--139, 2017.

\bibitem{meta2024llama3}
Meta AI,
``Llama 3: The most capable openly available LLM to date,''
Meta AI Technical Report, 2024.
\url{https://ai.meta.com/blog/meta-llama-3/}

\bibitem{mnih2015}
V.~Mnih, K.~Kavukcuoglu, D.~Silver, et~al.,
``Human-level control through deep reinforcement learning,''
\textit{Nature}, vol.~518, no.~7540, pp.~529--533, 2015.

\bibitem{chen2018}
X.~Chen, M.~Zhan, and H.~Zhou,
``Optimal evacuation route planning of urban populations under hazard scenarios using a hybrid simulation optimization approach,''
\textit{Safety Science}, vol.~107, pp.~190--201, 2018.

\bibitem{yao2023react}
S.~Yao, J.~Zhao, D.~Yu, N.~Du, I.~Shafran, K.~Narasimhan, and Y.~Cao,
``ReAct: Synergizing reasoning and acting in language models,''
in \textit{Proceedings of the 11th International Conference on Learning Representations (ICLR)}, 2023.

\bibitem{cho2014gru}
K.~Cho, B.~van Merrienboer, C.~Gulcehre, D.~Bahdanau, F.~Bougares, H.~Schwenk, and Y.~Bengio,
``Learning phrase representations using RNN encoder-decoder for statistical machine translation,''
in \textit{Proceedings of the 2014 Conference on Empirical Methods in Natural Language Processing (EMNLP)}, pp.~1724--1734, 2014.

\bibitem{bengio2009curriculum}
Y.~Bengio, J.~Louradour, R.~Collobert, and J.~Weston,
``Curriculum learning,''
in \textit{Proceedings of the 26th International Conference on Machine Learning (ICML)}, pp.~41--48, 2009.

\bibitem{osm}
OpenStreetMap contributors,
``OpenStreetMap,''
\url{https://www.openstreetmap.org}, 2024.

\end{thebibliography}

\end{document}
